\structure{ОСНОВНАЯ ЧАСТЬ}

В работе рассматривается геораспределённая система межкластерной
репликации, где один из них является активным, а другой пассивным.
Активный кластер принимает операции чтения и записи,
формирует упорядоченный журнал изменений
и асинхронно передаёт его в пассивный кластер.
Пассивный кластер последовательно применяет записи журнала,
поддерживая актуальную резервную копию состояния.

Внутри каждого кластера используется механизм консенсуса Raft~\cite{raft},
который обеспечивает согласованное применение операций
и устойчивость к отказу части узлов.
Переключение роли активного кластера выполняется
на основе механизма обнаружения отказов
и заданных тайм-аутов,
что позволяет обеспечивать отказоустойчивость системы.

Система должна обеспечивать следующие гарантии:
\begin{itemize}
\item целостность реплицируемых данных —
отсутствие подмены, потери и некорректного применения записей журнала;
\item доступность операций чтения и записи
при отказе части узлов и кратковременных сетевых сбоях;
\item ограниченное отставание репликации между активным и пассивным кластерами.
\end{itemize}

Реализация указанных гарантий осложняется
случайными сетевыми задержками,
возможными потерями сообщений
и инфраструктурными отказами узлов.
В связи с этим возникает необходимость
в математической постановке задачи,
позволяющей количественно описать критерии
целостности, доступности и согласованности данных
и сформулировать задачу выбора параметров системы,
обеспечивающих выполнение указанных требований.

\section{Характеристика организации}

Научно-учебный комплекс «Информатика и системы управления»~\cite{BMSTU_IU}
(далее — НУК ИУ) является самостоятельным структурным подразделением
МГТУ им. Н.Э. Баумана (далее — Университет).
Полное наименование НУК ИУ на английском языке —
Scientific Educational Complex «Informatics and Control Systems»
Bauman Moscow State Technical University.
Сокращённое наименование — SEC ICS BMSTU.

НУК ИУ возглавляет руководитель НУК ИУ,
доктор технических наук, профессор,
заведующий кафедрой ИУ6 («Компьютерные системы и сети»)
Пролетарский А.В.,
который непосредственно подчиняется
действующему ректору Университета
кандидату технических наук Гордину М.В.

Цель деятельности НУК ИУ —
обеспечение образовательного и научного процессов
в соответствии с Уставом Университета.

К основным задачам подразделения относятся:

\begin{itemize}
\item подготовка бакалавров, специалистов, магистров,
а также научных и научно-педагогических кадров
высшей квалификации в соответствии с государственными лицензиями,
выданными Университету;

\item выполнение фундаментальных, поисковых и прикладных
научных исследований, проведение опытно-конструкторских работ
и создание опытных образцов перспективной техники и технологий;

\item написание и издание учебников, учебных пособий и монографий;

\item развитие научных, педагогических и инженерных школ НУК ИУ;

\item разработка и внедрение прогрессивных форм, методов и средств
подготовки специалистов;

\item повышение квалификации и переподготовка научно-педагогических кадров
технических и экономических образовательных организаций,
а также специалистов предприятий с высшим образованием;

\item развитие материально-технической базы НУК ИУ;

\item организация научной работы студентов.
\end{itemize}
